% GENERATED FILE. Edit canonical lessons, then run npm run generate:books.
% canonical-source-hash: fnv1a64:8b21b21190df8508
% canonical-lessons: MR-C163-thand-karne, MR-C163-rikame-karne, MR-C163-vajan-karne, MR-C163-tapasne, MR-C163-sahi-karne

\chapter{To cool, To empty, To weigh, To check, To sign}
\label{ch:mr-to-cool-to-empty-to-weigh-to-check-to-sign}
\hlchaptermodality{163}

\begin{chapteropening}
\textbf{By the end of this chapter:} \emph{I can say to cool, to empty, to weigh, to check and to sign.}

Complete the last lesson of chapter 163: I can say to cool, to empty, to weigh, to check and to sign.
\end{chapteropening}

\section[\texorpdfstring{\mr{थंड} \mr{करणे} (thaṇḍ karṇe)}{थंड करणे (thaṇḍ karṇe)}]{\mr{थंड} \mr{करणे} (thaṇḍ karṇe) — to cool}

\label{lesson:MR-C163-thand-karne}



\begin{quote}
Before the new one: say the Marathi for to call over, then the Marathi for to thank.
\end{quote}

\subsection*{You'll want to know: \mr{थंड} \mr{करणे}}
\textbf{\mr{थंड} \mr{करणे}} — \emph{thaṇḍ karṇe} — \textquotedblleft{}to cool\textquotedblright{}.

\begin{grammarlens}[title={What you've built}]
One more word for this chapter.
\end{grammarlens}

\subsection*{Your turn}
\begin{itemize}
\raggedright
  \item \emph{Say it:} \emph{thaṇḍ karṇe}
  \item \emph{Say it:} \emph{thaṇḍ karṇe}, once more
  \item \emph{Say it:} say \emph{thaṇḍ karṇe}
  \item \emph{Recall:} say the Marathi for to call over, then the Marathi for to thank, then say \emph{thaṇḍ karṇe} again
\end{itemize}

\begin{tcolorbox}[breakable,colback=teal!4,colframe=teal!35!black,title={Before you move on}]
What does \mr{थंड} \mr{करणे} mean? (\textquotedblleft{}To cool\textquotedblright{}.) Say it once more.
\end{tcolorbox}
\section[\texorpdfstring{\mr{रिकामे} \mr{करणे} (rikāme karṇe)}{रिकामे करणे (rikāme karṇe)}]{\mr{रिकामे} \mr{करणे} (rikāme karṇe) — to empty}

\label{lesson:MR-C163-rikame-karne}



\begin{quote}
Before the new one: say the Marathi for to thank, then the Marathi for to cool.
\end{quote}

\subsection*{You'll want to know: \mr{रिकामे} \mr{करणे}}
\textbf{\mr{रिकामे} \mr{करणे}} — \emph{rikāme karṇe} — \textquotedblleft{}to empty\textquotedblright{}.

\begin{grammarlens}[title={What you've built}]
One more word for this chapter.
\end{grammarlens}

\subsection*{Your turn}
\begin{itemize}
\raggedright
  \item \emph{Say it:} \emph{rikāme karṇe}
  \item \emph{Say it:} \emph{rikāme karṇe}, once more
  \item \emph{Say it:} say \emph{rikāme karṇe}
  \item \emph{Recall:} say the Marathi for to thank, then the Marathi for to cool, then say \emph{rikāme karṇe} again
\end{itemize}

\begin{tcolorbox}[breakable,colback=teal!4,colframe=teal!35!black,title={Before you move on}]
What does \mr{रिकामे} \mr{करणे} mean? (\textquotedblleft{}To empty\textquotedblright{}.) Say it once more.
\end{tcolorbox}
\section[\texorpdfstring{\mr{वजन} \mr{करणे} (vajan karṇe)}{वजन करणे (vajan karṇe)}]{\mr{वजन} \mr{करणे} (vajan karṇe) — to weigh}

\label{lesson:MR-C163-vajan-karne}



\begin{quote}
Before the new one: say the Marathi for to cool, then the Marathi for to empty.
\end{quote}

\subsection*{You'll want to know: \mr{वजन} \mr{करणे}}
\textbf{\mr{वजन} \mr{करणे}} — \emph{vajan karṇe} — \textquotedblleft{}to weigh\textquotedblright{}.

\begin{grammarlens}[title={What you've built}]
One more word for this chapter.
\end{grammarlens}

\subsection*{Your turn}
\begin{itemize}
\raggedright
  \item \emph{Say it:} \emph{vajan karṇe}
  \item \emph{Say it:} \emph{vajan karṇe}, once more
  \item \emph{Say it:} say \emph{vajan karṇe}
  \item \emph{Recall:} say the Marathi for to cool, then the Marathi for to empty, then say \emph{vajan karṇe} again
\end{itemize}

\begin{tcolorbox}[breakable,colback=teal!4,colframe=teal!35!black,title={Before you move on}]
What does \mr{वजन} \mr{करणे} mean? (\textquotedblleft{}To weigh\textquotedblright{}.) Say it once more.
\end{tcolorbox}
\section[\texorpdfstring{\mr{तपासणे} (tapāsṇe)}{तपासणे (tapāsṇe)}]{\mr{तपासणे} (tapāsṇe) — to check}

\label{lesson:MR-C163-tapasne}



\begin{quote}
Before the new one: say the Marathi for to empty, then the Marathi for to weigh.
\end{quote}

\subsection*{You'll want to know: \mr{तपासणे}}
\textbf{\mr{तपासणे}} — \emph{tapāsṇe} — \textquotedblleft{}to check\textquotedblright{}.

\begin{grammarlens}[title={What you've built}]
One more word for this chapter.
\end{grammarlens}

\subsection*{Your turn}
\begin{itemize}
\raggedright
  \item \emph{Say it:} \emph{tapāsṇe}
  \item \emph{Say it:} \emph{tapāsṇe}, once more
  \item \emph{Say it:} say \emph{tapāsṇe}
  \item \emph{Recall:} say the Marathi for to empty, then the Marathi for to weigh, then say \emph{tapāsṇe} again
\end{itemize}

\begin{tcolorbox}[breakable,colback=teal!4,colframe=teal!35!black,title={Before you move on}]
What does \mr{तपासणे} mean? (\textquotedblleft{}To check\textquotedblright{}.) Say it once more.
\end{tcolorbox}
\section[\texorpdfstring{\mr{सही} \mr{करणे} (sahī karṇe)}{सही करणे (sahī karṇe)}]{\mr{सही} \mr{करणे} (sahī karṇe) — to sign}

\label{lesson:MR-C163-sahi-karne}



\begin{quote}
Before the new one: say the Marathi for to weigh, then the Marathi for to check.
\end{quote}

\subsection*{You'll want to know: \mr{सही} \mr{करणे}}
\textbf{\mr{सही} \mr{करणे}} — \emph{sahī karṇe} — \textquotedblleft{}to sign\textquotedblright{}.

\begin{grammarlens}[title={What you've built}]
One more word for this chapter.
\end{grammarlens}

\subsection*{Your turn}
\begin{itemize}
\raggedright
  \item \emph{Say it:} \emph{sahī karṇe}
  \item \emph{Say it:} \emph{sahī karṇe}, once more
  \item \emph{Say it:} say \emph{sahī karṇe}
  \item \emph{Recall:} say the Marathi for to weigh, then the Marathi for to check, then say \emph{sahī karṇe} again
\end{itemize}

\begin{tcolorbox}[breakable,colback=teal!4,colframe=teal!35!black,title={Before you move on}]
What does \mr{सही} \mr{करणे} mean? (\textquotedblleft{}To sign\textquotedblright{}.) Say it once more.
\end{tcolorbox}
